% Part: incompleteness
% Chapter: incompleteness-provability
% Section: introduction

\documentclass[../../../include/open-logic-section]{subfiles}

\begin{document}

\olfileid{inc}{inp}{int}

\olsection{ভূমিকা}

হিলবার্ট মনে করতেন, স্বাভাবিক সংখ্যার মতো কোনো গাণিতিক গঠনের জন্য
স্বতঃসিদ্ধের ব্যবস্থা তখনই যথেষ্ট, যখন সেটি থেকে ওই গঠন সম্বন্ধে সব সত্য
উক্তি নিষ্পাদন করা যায়। পরে বিধিবদ্ধ নিষ্পাদন-ব্যবস্থায় তাঁর আগ্রহের
সঙ্গে এই ধারণা মিলিয়ে বোঝা যায়, তাঁর মতে আমাদের নিশ্চিত করা উচিত যে,
ধরা যাক, স্বাভাবিক সংখ্যা নিয়ে বিচার করতে ব্যবহৃত বিধিবদ্ধ ব্যবস্থাগুলি
শুধু সঙ্গতিপূর্ণ নয়, \emph{সম্পূর্ণ}-ও; অর্থাৎ তার ভাষার প্রতিটি উক্তি
হয় !!{derivable}, নয়তো তার নঞর্থক রূপটি তা-ই। গ্যোডেলের প্রথম
অসম্পূর্ণতা উপপাদ্য দেখায় যে এমন কোনো স্বতঃসিদ্ধ-ব্যবস্থা নেই: পাটিগণিতের
কোনো সম্পূর্ণ, সঙ্গতিপূর্ণ, !!{axiomatizable} বিধিবদ্ধ ব্যবস্থা নেই। বস্তুত,
কোনো ``যথেষ্ট শক্তিশালী,'' সঙ্গতিপূর্ণ, !!{axiomatizable} গাণিতিক তত্ত্বই
সম্পূর্ণ নয়।

হিলবার্টের আরও গুরুত্বপূর্ণ লক্ষ্য—আধুনিক (``ধ্রুপদি'') গণিতের যৌক্তিকতা
প্রতিষ্ঠার কর্মসূচির কেন্দ্রবিন্দু—ছিল ধ্রুপদি যুক্তিবিচারকে প্রকাশকারী
বিধিবদ্ধ ব্যবস্থাগুলির সঙ্গতির সসীমতাবাদী প্রমাণ খুঁজে পাওয়া। তাই
হিলবার্টের কর্মসূচির পক্ষে গ্যোডেলের দ্বিতীয় অসম্পূর্ণতা উপপাদ্য আরও বড়
আঘাত ছিল। প্রথমটির মতো দ্বিতীয় অসম্পূর্ণতা উপপাদ্যও অস্পষ্ট ভাষায় বলা
যায়। মোটামুটি, কোনো যথেষ্ট শক্তিশালী পাটিগাণিতিক তত্ত্ব নিজের সঙ্গতি
প্রমাণ করতে পারে না। এখানে ``যথেষ্ট শক্তিশালী'' বলতে আমাদের
$\Th{Q}$-এর চেয়ে সামান্য বেশি কিছু বোঝাতে হবে।

অসম্পূর্ণতা উপপাদ্যের গ্যোডেলীয় মূল প্রমাণের ধারণাটি এপিমেনিদেসের
কূটাভাসে পাওয়া যায়। ক্রিটের অধিবাসী এপিমেনিদেস বলেছিলেন যে ক্রিটের সব
অধিবাসী মিথ্যাবাদী; কূটাভাসটির আরও সরাসরি রূপ হলো ``এই বাক্যটি মিথ্যা।''
মূলত সত্যের জায়গায় !!{derivability} বসিয়ে গ্যোডেল এমন !!a{sentence}-কে
বিধিবদ্ধ করতে পেরেছিলেন, যা ঘুরিয়ে বলে যে সেটি নিজেই !!{derivable} নয়।
ওই !!{sentence}টি !!{derivable} হলে তত্ত্বটি অসঙ্গতিপূর্ণ হতো। গ্যোডেল আরও
দেখিয়েছিলেন, তাঁর বিবেচিত স্বতঃসিদ্ধ-ব্যবস্থা থেকে ওই !!{sentence}-এর
নঞর্থক রূপটিও !!{derivable} নয়। (এই দ্বিতীয় অংশে গ্যোডেলকে ধরে নিতে
হয়েছিল যে তত্ত্ব~$\Th{T}$ ``$\omega$-সঙ্গতিপূর্ণ।'' $\omega$-সঙ্গতি
সঙ্গতির সঙ্গে সম্পর্কিত, কিন্তু তার চেয়ে শক্তিশালী ধর্ম।\footnote{অর্থাৎ,
প্রতিটি $\omega$-সঙ্গতিপূর্ণ তত্ত্ব সঙ্গতিপূর্ণ, কিন্তু বিপরীতটি সত্য নয়।}
গ্যোডেলের কয়েক বছর পরে রসার দেখান যে $\Th{T}$-এর সাধারণ সঙ্গতিই যথেষ্ট।)

প্রথম চ্যালেঞ্জ হলো, নিজেকেই উল্লেখ করে এমন !!a{sentence} কীভাবে নির্মাণ
করা যায় তা বোঝা। $\Th{Q}$-এর ভাষার প্রতিটি !!{formula}~$!A$-এর জন্য
$\gn{!A}$ দিয়ে $\Gn{!A}$-এর সংশ্লিষ্ট সংখ্যাপদটি বোঝাই। এর অর্থটি খেয়াল
করো: $!A$ হলো $\Th{Q}$-এর ভাষার !!a{formula}, $\Gn{!A}$ একটি স্বাভাবিক
সংখ্যা, আর $\gn{!A}$ হলো $\Th{Q}$-এর ভাষার একটি \emph{পদ}। অতএব
$\Th{Q}$-এর ভাষার প্রতিটি !!{formula}~$!A$-এর একটি \emph{নাম},
$\gn{!A}$, আছে, যা $\Th{Q}$-এর ভাষার একটি পদ; এই ব্যবস্থা $\Th{Q}$-এর
ভাষার !!{formula}s-কে অন্য !!{formula}s সম্বন্ধে কিছু ``বলতে'' পারার
ধারণাগত পরিকাঠামো দেয়। নিচের লেমাটি স্থির-বিন্দু লেমা নামে পরিচিত।

\begin{lem}
ধরি $\Th{T}$ হলো $\Th{Q}$-কে প্রসারিত করে এমন যেকোনো তত্ত্ব, এবং
$!B(x)$ কেবল $x$ চলটি মুক্ত এমন যেকোনো !!{formula}। তবে এমন
!!a{sentence}~$!A$ আছে, যাতে $\Th{T} \Proves!A \liff !B(\gn{!A})$।
\end{lem}

লেমাটি বলে, যেকোনো ধর্ম $!B(x)$ দেওয়া থাকলে এমন !!a{sentence}~$!A$
আছে যা বলে ``$!B(x)$ আমার ক্ষেত্রে সত্য,'' এবং $\Th{T}$ এটি ``জানে।''

এমন !!a{sentence} আমরা কীভাবে নির্মাণ করতে পারি? কুয়াইনের দেওয়া
এপিমেনিদেস-কূটাভাসের নিচের রূপটি বিবেচনা করো:
\begin{quote}
``নিজের উদ্ধৃত রূপটি আগে বসালে মিথ্যা দেয়''-এর আগে তার উদ্ধৃত রূপটি
বসালে মিথ্যা দেয়।
\end{quote}
এই বাক্যটি সরাসরি নিজেকে উল্লেখ করে না। এটি শুধু উদ্ধৃতিচিহ্নের মধ্যের
সংকেতবিন্যাসগত বস্তুগুলি সম্বন্ধে একটি কথা বলে; এই দিক থেকে এটি নিচের
বাক্যগুলির সমপর্যায়ের:
\begin{enumerate}
\item ``রবার্ট'' একটি সুন্দর নাম।
\item ``আমি দৌড়েছিলাম।'' একটি ছোট বাক্য।
\item ``তিনটি শব্দ আছে''-তে তিনটি শব্দ আছে।
\end{enumerate}
কিন্তু ``নিজের উদ্ধৃত রূপটি আগে বসালে মিথ্যা দেয়'' পদবন্ধটির আগে তারই
একটি উদ্ধৃত রূপ বসালে কী হয়? তখন আবার মূল বাক্যটিই পাওয়া যায়!{} সংক্ষেপে,
বাক্যটি বলে যে সেটি মিথ্যা।

\end{document}
